% BEGIN VOL08-EXPANSION VIII21
\section{Supplementary worked examples}
\label{sec:viii21-pedagogy}

\begin{example}[Euler characteristic of a sphere]\label{ex:viii21-ped-01}
The CW structure on $S^n$ with one $0$-cell and one $n$-cell gives
\[
\chi(S^n)=1+(-1)^n.
\]
The same number is obtained from Betti numbers, illustrating the Euler--Poincare identity.
\end{example}

\begin{example}[Euler characteristic of a closed orientable surface]\label{ex:viii21-ped-02}
A genus-$g$ surface has a CW structure with one vertex, $2g$ one-cells, and one two-cell. Therefore
\[
\chi(\Sigma_g)=1-2g+1=2-2g.
\]
\end{example}

\begin{example}[Inclusion--exclusion]\label{ex:viii21-ped-03}
For finite CW subcomplexes $A,B$ with $X=A\cup B$,
\[
\chi(X)=\chi(A)+\chi(B)-\chi(A\cap B).
\]
This is the Euler-characteristic shadow of the Mayer--Vietoris exact sequence.
\end{example}

\section{Graded supplementary exercises}
The following exercises progress from direct computation and construction to proof, hypothesis testing, application, and synthesis.

\subsection*{Standard computations and constructions}

\begin{exercise}[Circle]\label{exr:viii21-ped-01}
Compute $\chi(S^1)$.
\end{exercise}
\begin{hint}
Use one $0$-cell and one $1$-cell.
\end{hint}
\begin{solution}
$\chi(S^1)=1-1=0$.
\end{solution}

\begin{exercise}[Two-sphere]\label{exr:viii21-ped-02}
Compute $\chi(S^2)$.
\end{exercise}
\begin{hint}
Use one $0$-cell and one $2$-cell.
\end{hint}
\begin{solution}
$\chi(S^2)=2$.
\end{solution}

\begin{exercise}[Torus]\label{exr:viii21-ped-03}
Compute $\chi(T^2)$ from its standard CW structure.
\end{exercise}
\begin{hint}
Use cell counts $1,2,1$.
\end{hint}
\begin{solution}
$\chi(T^2)=1-2+1=0$.
\end{solution}

\begin{exercise}[Wedge of circles]\label{exr:viii21-ped-04}
Compute $\chi(\bigvee_{i=1}^r S^1)$.
\end{exercise}
\begin{hint}
Use one vertex and $r$ edges.
\end{hint}
\begin{solution}
$\chi=1-r$.
\end{solution}

\begin{exercise}[Projective plane]\label{exr:viii21-ped-05}
Compute $\chi(\mathbb{RP}^2)$ from its one-cell-per-dimension CW structure.
\end{exercise}
\begin{hint}
Count dimensions $0,1,2$.
\end{hint}
\begin{solution}
$\chi(\mathbb{RP}^2)=1-1+1=1$.
\end{solution}

\subsection*{Proofs}

\begin{exercise}[Euler--Poincare]\label{exr:viii21-ped-06}
Prove for a finite free chain complex that $\sum(-1)^n\operatorname{rank}C_n=\sum(-1)^n\operatorname{rank}H_n$.
\end{exercise}
\begin{hint}
Split ranks into boundaries, homology, and shifted boundaries.
\end{hint}
\begin{solution}
Since $\operatorname{rank}C_n=\operatorname{rank}B_n+\operatorname{rank}H_n+\operatorname{rank}B_{n-1}$, the boundary ranks cancel pairwise in the alternating sum.
\end{solution}

\begin{exercise}[Homotopy invariance of chi]\label{exr:viii21-ped-07}
Prove $\chi$ is a homotopy invariant for finite CW complexes.
\end{exercise}
\begin{hint}
Use homology ranks.
\end{hint}
\begin{solution}
Homotopy-equivalent finite CW complexes have isomorphic homology groups, so the alternating sum of Betti numbers agrees.
\end{solution}

\begin{exercise}[Wedge formula]\label{exr:viii21-ped-08}
Prove $\chi(X\vee Y)=\chi(X)+\chi(Y)-1$ for connected finite CW complexes.
\end{exercise}
\begin{hint}
Apply inclusion--exclusion at the wedge point.
\end{hint}
\begin{solution}
The intersection is a point of Euler characteristic $1$, yielding the formula.
\end{solution}

\begin{exercise}[Product formula]\label{exr:viii21-ped-09}
Prove $\chi(X\times Y)=\chi(X)\chi(Y)$ for finite CW complexes.
\end{exercise}
\begin{hint}
Product cells have dimensions that add.
\end{hint}
\begin{solution}
Summing $(-1)^{p+q}c_p(X)c_q(Y)$ factors into the product of the two alternating sums.
\end{solution}

\subsection*{Counterexamples and hypothesis tests}

\begin{exercise}[Cell decomposition dependence]\label{exr:viii21-ped-10}
Test: $\chi(X)$ depends on the chosen finite CW decomposition.
\end{exercise}
\begin{hint}
Compare the homology formula.
\end{hint}
\begin{solution}
False. Different finite CW structures give the same alternating sum because both equal the homological Euler characteristic.
\end{solution}

\begin{exercise}[Zero chi means torus]\label{exr:viii21-ped-11}
Test: a connected closed surface with $\chi=0$ must be a torus.
\end{exercise}
\begin{hint}
Remember nonorientable surfaces.
\end{hint}
\begin{solution}
False without orientability; the Klein bottle also has Euler characteristic zero.
\end{solution}

\begin{exercise}[Chi determines homology]\label{exr:viii21-ped-12}
Test: equal Euler characteristics imply isomorphic homology groups.
\end{exercise}
\begin{hint}
An alternating sum loses information.
\end{hint}
\begin{solution}
False. Many different Betti-number patterns can have the same alternating sum.
\end{solution}

\subsection*{Applications and investigations}

\begin{exercise}[Mesh sanity check]\label{exr:viii21-ped-13}
Explain the formula $\chi=V-E+F$ for a finite triangulated surface.
\end{exercise}
\begin{hint}
Use simplices as cells.
\end{hint}
\begin{solution}
Vertices, edges, and faces are the $0$-, $1$-, and $2$-cells, so the cellular alternating sum is $V-E+F$.
\end{solution}

\begin{exercise}[Genus recovery]\label{exr:viii21-ped-14}
For a connected closed orientable surface, recover genus from $\chi$.
\end{exercise}
\begin{hint}
Use $\chi=2-2g$.
\end{hint}
\begin{solution}
$g=1-\chi/2$.
\end{solution}

\subsection*{Challenge problems}

\begin{exercise}[Connected sum]\label{exr:viii21-ped-15}
Derive $\chi(M\#N)=\chi(M)+\chi(N)-2$ for closed surfaces.
\end{exercise}
\begin{hint}
Remove a disk from each and glue along $S^1$.
\end{hint}
\begin{solution}
Removing each open disk lowers $\chi$ by $1$; gluing along a circle contributes no correction since $\chi(S^1)=0$, giving the formula.
\end{solution}

\begin{exercise}[Euler synthesis]\label{exr:viii21-ped-16}
Explain three independent routes to $\chi$ and when each is useful.
\end{exercise}
\begin{hint}
Think cells, homology, and inclusion--exclusion.
\end{hint}
\begin{solution}
One may alternate cell counts for a CW model, alternate Betti numbers after homology computation, or decompose the space and use additivity; agreement is a strong consistency check.
\end{solution}

% END VOL08-EXPANSION VIII21
